%% ma: L12-B19-0011
%% lop: 12
%% chuong: 6
%% bai: 19
%% dang: 12.19.1
%% loai: TLN
%% muc: TH
%% dap_an: "0,29"
%% nguon: "(NBV)-ĐỀ SỐ 6-MỖI NGÀY 1 ĐỀ THI 2025.pdf (D06, MỖI NGÀY 1 ĐỀ - Đề số 6), Phần III Câu 6"
%% kien_thuc: "Bài 19 (công thức xác suất toàn phần)"
%% trang_thai: chinh-thuc
%% ghi_chu: "Trùng ý tưởng với dạng 12.19.1 (công thức xác suất toàn phần trong khảo sát) nên nhập cùng dạng; đáp số là xác suất toàn phần, giữ ở mức TH. Bỏ ảnh minh họa người hút thuốc (ảnh đời sống, không chứa dữ kiện toán). Đáp án gốc 0,29 đúng"
\begin{ex}
Tỉ lệ người dân nghiện thuốc là $30\%$, biết rằng tỉ lệ người viêm phổi trong số người nghiện thuốc lá là $50\%$, còn tỉ lệ người viêm phổi trong số người không hút thuốc là $20\%$. Chọn ngẫu nhiên $1$ người. Tìm xác suất để người đó bị viêm phổi.
\shortans{0,29}
\loigiai{Gọi $A$ là biến cố ``người được chọn bị viêm phổi'', $B$ là biến cố ``người được chọn nghiện thuốc''. Ta có $P(B)=0{,}3$, $P(\overline B)=0{,}7$, $P(A\mid B)=0{,}5$, $P(A\mid\overline B)=0{,}2$. Theo công thức xác suất toàn phần:
\[P(A)=P(B)P(A\mid B)+P(\overline B)P(A\mid\overline B)=0{,}3\cdot0{,}5+0{,}7\cdot0{,}2=0{,}29.\]}
\end{ex}
